
Debiasing methods such as Just Train Twice (JTT) rely on the assumption that spurious features---shortcuts like color or background cues---are learned earlier than core features during gradient descent training, but this temporal ordering has never been directly measured at the gradient level. This work validates this assumption for the first time via ablation training and per-epoch gradient tracking on CMNIST. Spurious features (color) converge 4 epochs earlier than core features (shape) during gradient descent ($E_s=13$ vs $E_c=17, \Delta=4$ epochs, single-dataset proof-of-concept with seed 0; full 10-seed statistical validation in progress). Early convolutional layers exhibit significantly higher neuron-spurious correlation than late layers (mean $\rho_j = 0.004$ vs $0.001, p=0.0028, t=2.78$, Cohen's $d=0.25$), confirming that temporal ordering arises from architectural feature hierarchy where simpler low-level features stabilize before high-level semantics. Spurious-trained networks show 51\% lower forgetting rate (2.35 vs 4.82 events/sample), confirming stability advantages. However, gradient-aware training using CMNIST-derived neuron correlations failed on Waterbirds (39\% worst-group accuracy vs 86\% JTT target), revealing that neuron correlations are dataset-specific and cannot transfer across spurious feature types. Results validate the mechanistic foundation underlying JTT/LfF reweighting methods while identifying cross-dataset neuron transfer as a fundamental constraint for neuron-level interventions. Generalization beyond color-based spurious correlations to background (Waterbirds), attribute (CelebA), and context (NICO++) is future work.
